\section{Results}

\subsection{RQ1: Real Execution}
\textbf{Answer to RQ1.} Yes for real execution and reproducibility, within the scope of the nominal benchmark-aligned paper campaign. The canonical nominal run executed 28 circuits, all in \texttt{REAL} mode, with no mock result retained in the scientific evidence, 28 successful executions, and 28 scientifically eligible outputs.

The direct evidence comes from the canonical paper campaign \texttt{results/paper\_campaign\_summary.json} and \texttt{results/paper\_campaign\_summary.csv}, run id \texttt{20260711\_094959}. All 28 rows have \texttt{simulation\_mode = REAL}, \texttt{execution\_status = SUCCESS}, and \texttt{paper\_eligible = True}. The same campaign covers multiple analysis families through the extracted nominal metrics: operating-point and DC quantities such as operating point and quiescent current, AC quantities such as \texttt{dc\_gain\_db} and cutoff frequency, transient quantities such as propagation delay, slew rate, settling time, startup amplitude, and oscillator frequency, and spectral quantities such as THD and fundamental frequency. These results establish that Spec2Testbench can produce real and traceable verification evidence across the benchmark-aligned corpus.

The interpretation is limited to execution and evidence production. Successful real execution does not by itself establish that a circuit is specification-compliant, and the next subsection shows why this distinction is necessary.

\subsection{RQ2: Simulability Versus Compliance}
\textbf{Answer to RQ2.} Yes. The canonical nominal campaign shows that Spec2Testbench distinguishes simulability from compliance: 28 circuits are simulable in the nominal campaign, but only 27 are compliant and one is noncompliant.

The nominal result is directly visible in \path{results/paper_campaign_summary.csv} and \path{results/simulability_vs_compliance.csv}. The campaign contains 27 \texttt{SIMULABLE\_COMPLIANT} cases and one \texttt{SIMULABLE\_NONCOMPLIANT} case. The unique noncompliant nominal case is \texttt{p04\_amplifier}. Its per-case artifact \path{artifacts/paper_campaign/20260711_094959/p04_amplifier/report.json} shows \texttt{execution\_status = SUCCESS}, \texttt{simulation\_mode = REAL}, and \texttt{scientific\_category = SIMULABLE\_NONCOMPLIANT}. The metric record in \path{results/paper_metric_results.csv} shows that \texttt{dc\_gain\_db} was extracted successfully with measured value \(-160.0000000868589\) dB against the requirement \(\ge 0.0\) dB, yielding a noncompliant verdict.

This case provides the direct counterexample needed by the paper:
\[
\mathrm{SIMULABLE} \neq \mathrm{COMPLIANT}
\]
Spec2Testbench therefore does not collapse successful ngspice termination into specification satisfaction. The result is scientifically useful precisely because it preserves the difference between executable evidence and compliant behavior.

\subsection{RQ3: Controlled Violations}
\textbf{Answer to RQ3.} The canonical evidence does not support a claim of perfect controlled-violation detection across all executed datasets. Instead, it supports two distinct statements: the frozen pilot V3 achieved perfect separation within its limited 16-case scope, whereas the later expanded controlled-violation campaign retained one false accept among two effective violations.

The retained frozen pilot evidence is \texttt{results/frozen\_pilot\_metrics\_v3.json}. This dataset records 16 cases with \texttt{TRUE\_ACCEPT = 8}, \texttt{TRUE\_DETECTION = 8}, \texttt{FALSE\_ACCEPT = 0}, \texttt{FALSE\_REJECT = 0}, and \texttt{UNEVALUATED = 0}. It also includes two WRDATA extension cases and is the canonical source for the small-scope backend-validated pilot. The later expanded campaign is \texttt{results/final\_controlled\_summary.json}, dated by run id \texttt{20260712\_195226}. That campaign generated 30 variants, but only 2 were effective controlled violations and 28 were ineffective mutations. Within those 2 effective violations, one was detected and one was missed, yielding one false pass and zero unevaluated effective cases. This later campaign is therefore the canonical source for the broader sensitivity analysis, not for a perfect-detection claim.

The manuscript should not silently merge these datasets. The frozen pilot V3 is retained as the canonical small-scope campaign for a clean confusion matrix because it is internally consistent, backend-validated, and machine-readable. The contradictory later campaign must be reported explicitly as a limitation or sensitivity result: it shows that the broader mutation population was far harder than the frozen pilot and that perfect controlled-violation detection did not generalize to the expanded effective set.

The interpretation is therefore two-level. In the frozen pilot V3, Spec2Testbench achieved perfect discrimination on the retained 16-case pilot. In the later expanded mutation campaign, the framework still separated effective from ineffective mutation accounting, but one effective violation remained falsely accepted. This is a limitation of the broader controlled-violation evidence, not a reason to overwrite the frozen pilot counts.

\subsection{RQ4: Baseline Versus LLM}
\textbf{Answer to RQ4.} No quantitative answer can be given from the current canonical evidence because the required ablation results are not yet available.

The canonical ablation artifact \texttt{results/final\_ablation\_summary.json} records \texttt{status = PARTIAL}, \texttt{A2 = NOT\_EXECUTED}, \texttt{A4 = NOT\_EXECUTED}, and \texttt{llm\_included = false}. The canonical nominal paper campaign configuration in \texttt{configs/paper\_experiment.yaml} also sets \texttt{llm.enabled: false}. As a result, the manuscript cannot report valid generation rate, simulation rate, metric coverage, correct verdict rate, false accept rate, cost, or run-to-run variability for an executed LLM-versus-deterministic comparison. Table~\ref{tab:planned_llm_ablation} is therefore intentionally left as a pending-results table rather than a quantitative result table.

The interpretation is straightforward: the LLM layer remains a planned experimental factor rather than a completed results block in the current manuscript evidence.

\subsection{Secondary Software Results}
\textbf{Answer to the secondary software question.} The software evidence supports backend functionality and reproducible execution infrastructure, but it is not used here as direct proof of analog specification compliance.

The canonical software-test artifact \texttt{results/final\_test\_results.json} records 66 unique tests with 66 passes, 0 failures, 0 skips, and 1 warning in both normal mode and PySpice-disabled mode. The backend evidence separately confirms that \texttt{NGSPICE\_MEASURE} and \texttt{NGSPICE\_WRDATA} both operate on canonical evidence, and \texttt{results/wrdata\_independent\_comparisons.csv} shows agreement in both WRDATA validation rows. These software and backend results establish that the execution and extraction machinery functions with and without PySpice, but they are presented here only as supporting infrastructure evidence. The analog-compliance claims in this section remain tied to the campaign artifacts and per-case metric records rather than to the software test suite.
